\section{AMC 12 12P 2002}
        \begin{enumerate}
        	\item (\textit{AMC 12 12P 2002}) Which of the following numbers is a perfect square?


 $\text{(A) }4^4 5^5 6^6 \qquad \text{(B) }4^4 5^6 6^5 \qquad \text{(C) }4^5 5^4 6^6 \qquad \text{(D) }4^6 5^4 6^5 \qquad \text{(E) }4^6 5^5 6^4$

 \textbf{Answer: }C



\textbf{Solution 1}

For a positive integer to be a perfect square, all the primes in its prime factorization must have an even exponent. With a quick glance at the answer choices, we can eliminate options

 $\textbf{(A)}$ because $5^5$ is an odd power

 $\textbf{(B)}$ because $6^5 = 2^5 \cdot 3^5$ and $3^5$ is an odd power

 $\textbf{(D)}$ because $6^5 = 2^5 \cdot 3^5$ and $3^5$ is an odd power, and

 $\textbf{(E)}$ because $5^5$ is an odd power.

 This leaves option $\textbf{(C)},$ in which $4^5=(2^{2})^{5}=2^{10}$, and since $10, 4,$ and $6$ are all even, $\textbf{(C)}$ is a perfect square. Thus, our answer is $\boxed{\textbf{(C) } 4^5 5^4 6^6}.$



	\item (\textit{AMC 12 12P 2002}) The function $f$ is given by the table


 
 \[\begin{tabular}{|c||c|c|c|c|c|}  \hline   x & 1 & 2 & 3 & 4 & 5 \\   \hline  f(x) & 4 & 1 & 3 & 5 & 2 \\  \hline \end{tabular}\]


 If $u_0=4$ and $u_{n+1} = f(u_n)$ for $n \ge 0$, find $u_{2002}$
$\text{(A) }1 \qquad \text{(B) }2 \qquad \text{(C) }3 \qquad \text{(D) }4 \qquad \text{(E) }5$

 \textbf{Answer: }B



\textbf{Solution 1}

We can guess that the series given by the problem is periodic in some way. Starting off, $u_0=4$ is given. $u_1=u_{0+1}=f(u_0)=f(4)=5,$ so $u_1=5.$ $u_2=u_{1+1}=f(u_1)=f(5)=2,$ so $u_2=2.$ $u_3=u_{2+1}=f(u_2)=f(2)=1,$ so $u_3=1.$ $u_4=u_{3+1}=f(u_3)=f(1)=4,$ so $u_4=4.$ Plugging in $4$ will give us $5$ as found before, and plugging in $5$ will give $2$ and so on. This means that our original guess of the series being periodic was correct. Summing up our findings in a nice table,

 
 \[\begin{tabular}{|c||c|c|c|c|c|c|}  \hline   n & 0 & 1 & 2 & 3 & 4 & ...\\   \hline  un & 4 & 5 & 2 & 1 & 4 & ...\\  \hline \end{tabular}\]

 in which the next $u_n$ is found by simply plugging in the number from the last box into $f(x).$ The function is periodic every $4$ terms. $2002 \equiv 2\pmod{4}$, and counting $4$ starting from $u_1$ will give us our answer of $\boxed{\textbf{(B) } 2}$.

 


	\item (\textit{AMC 12 12P 2002}) The dimensions of a rectangular box in inches are all positive integers and the volume of the box is $2002$ $\text{in}^3$. Find the minimum possible sum of the three dimensions.


 $\text{(A) }36 \qquad \text{(B) }38  \qquad \text{(C) }42 \qquad \text{(D) }44 \qquad \text{(E) }92$

 \textbf{Answer: }B



\textbf{Solution 1}

Given an arbitrary product and an arbitrary amount of positive integers to multiply to get that product, make the terms as close to each other as possible to minimize the sum. (This is a corollary that follows from the $AM-GM$ proof.) 

 A good rule of thumb is to memorize the prime factorization of the AMC year that you are doing, which is $2002= 2 \cdot 7 \cdot 11 \cdot 13$. The three terms that are closest to each other that multiply to $2002$ are $11, 13,$ and $14$, so our answer is $11+13+14=\boxed{\textbf{(B) } 38}$.

 


	\item (\textit{AMC 12 12P 2002}) Let $a$ and $b$ be distinct real numbers for which

 \[\frac{a}{b} + \frac{a+10b}{b+10a} = 2.\]


 Find $\frac{a}{b}$.


 $\text{(A) }0.4 \qquad \text{(B) }0.5  \qquad \text{(C) }0.6 \qquad \text{(D) }0.7 \qquad \text{(E) }0.8$

 \textbf{Answer: }E



\textbf{Solution 1}

For sake of speed, WLOG, let $b=1$. This means that the ratio $\frac{a}{b}$ will simply be $a$ because $\frac{a}{b}=\frac{a}{1}=a.$ Solving for $a$ with some very simple algebra gives us a quadratic which is $5a^2 -9a +4=0.$ Factoring the quadratic gives us $(5a-4)(a-1)=0$. Therefore, $a=1$ or $a=\frac{4}{5}=0.8.$ However, since $a$ and $b$ must be distinct, $a$ cannot be $1$ so the latter option is correct, giving us our answer of $\boxed{\textbf{(E) } 0.8}$ for the AMC 12P and $\boxed{\textbf{(C) } 0.8}$ for the AMC 10P.

 



\textbf{Solution 2}

The only tricky part about this equation is the fact that the left-hand side has fractions. Multiplying both sides by $b(b+10a)$ gives us $2ab+10a^2+10b^2=2b^2+20ab.$ Moving everything to the left-hand side and dividing by $2$ gives $5a^2-4b^2 -9ab,$ which factors into $(5a-4b)(a-b)=0.$ Because $a \neq b, 5a=4b \implies \frac{a}{b}=0.8$ giving us our answer of $\boxed{\textbf{(E) } 0.8}$ for the AMC 12P and $\boxed{\textbf{(C) } 0.8}$ for the AMC 10P.

 \textbf{Note} For some unknown reason, the answer choices for the \href{https://artofproblemsolving.com/wiki/index.php?title=2002\_AMC\_10P}{2002 AMC 10P} are different from the answer choices for the \href{https://artofproblemsolving.com/wiki/index.php?title=2002\_AMC\_12P}{2002 AMC 12P}, even though the question is exactly the same. Indeed, $\boxed{\textbf{(C) }0.8}$ is the correct answer choice for the 2002 AMC 10P.

 


	\item (\textit{AMC 12 12P 2002}) For how many positive integers $m$ is 

 \[\frac{2002}{m^2 -2}\]


 a positive integer?


 $\text{(A) one} \qquad \text{(B) two} \qquad \text{(C) three} \qquad \text{(D) four} \qquad \text{(E) more than four}$

 \textbf{Answer: }C



\textbf{Solution}

For $\frac{2002}{m^2 -2}$ to be an integer, $2002$ must be divisible by $m^2-2.$ Again, memorizing the prime factorization of $2002$ is helpful. $2002 = 2 \cdot 7 \cdot 11 \cdot 13$, so its factors are $1, 2, 7, 11, 13, 14, 22, 26, 77, 91, 143, 154, 182, 286, 1001$, and $2002$. 

 From this adding $2$ to the items in our list and checking if they are perfect squares will suffice. $4, 9,$ and $16$ end up being our perfect squares, giving us an answer of $\boxed{\textbf{(C) } \text{three}}.$



	\item (\textit{AMC 12 12P 2002}) Participation in the local soccer league this year is $10\%$ higher than last year. The number of males increased by $5\%$ and the number of females increased by $20\%$. What fraction of the soccer league is now female?


 $\text{(A) }\frac{1}{3} \qquad \text{(B) }\frac{4}{11} \qquad \text{(C) }\frac{2}{5} \qquad \text{(D) }\frac{4}{9} \qquad \text{(E) }\frac{1}{2}$

 \textbf{Answer: }B



\textbf{Solution 1}

Let the amount of soccer players last year be $x$, the number of male players last year to be $m$, and the number of females players last year to be $f.$ We want to find $\frac{1.2f}{1.1x},$ since that's the fraction of female players now. From the problem, we are given 

 \begin{align*}x&=m+f \\1.1x&=1.05m+1.2f \\\end{align*}


 Eliminating $m$ and solving for $\frac{1.2f}{1.1x}$ gives us our answer of $\boxed{\textbf{(B) } \frac {4}{11}}.$

 Also, you know the number has to be divisible by 11, and there’s only one answer that fits!:)

 


	\item (\textit{AMC 12 12P 2002}) How many three-digit numbers have at least one $2$ and at least one $3$?


 $\text{(A) }52 \qquad \text{(B) }54  \qquad \text{(C) }56 \qquad \text{(D) }58 \qquad \text{(E) }60$

 \textbf{Answer: }A



\textbf{Solution 1}

We can do this problem with some simple case work.

 Case 1: The hundreds place is not $2$ or $3.$This means that the tens place and ones place must be $2$ and $3$ respectively or $3$ and $2$ respectively. This case covers $1, 4, 5, 6, 7, 8,$ and $9,$ so it gives us $2 \cdot 7 = 14$ cases.

 Case 2: The hundreds place is $2.$This means that $3$ must be in the tens place or ones place. Starting with cases in which the tens place is not $3$, we get $203, 213, 223, 243, 253, 263, 273, 283,$ and $293.$  With cases in which the tens place is $3$, we have $230-239$, or $10$ more cases. This gives us $9 + 10=19$ cases.

 Case 3: The hundreds place is $3.$This case is almost identical to the second case, just swap the $2$s with $3$s and $3$s with $2$s in the reasoning and it's the same, giving us an additional $19$ cases.

 Adding up all of these cases gives $14+19+19=52$ cases, or $\boxed{\textbf{(A) }52}.$




\textbf{Solution 2 (PIE)}

We can solve this problem with the principle of inclusion-exclusion. Let $\#(A)$ be the number of three-digit numbers without the digit $2$, let $\#(B)$ be the number of three-digit numbers without the digit $3$, and let $\#(A \cap B)$ be the number of three-digit numbers without digits $2$ or $3$.

 Notice that $\#(A)=\#(B)$ by symmetry. We use constructive counting to find that 
 \[\#(A) = 8 \cdot 9 \cdot 9 = 648\] so $\#(A) = \#(B) = 648$. Similarly, 
 \[\#(A \cap B) = 7 \cdot 8 \cdot 8 = 448.\]

 Hence, by PIE, 
 \[\#(A \cup B) = \#(A) + \#(B) - \#(A \cap B) = 648 + 648 - 448 = 848\] so the count of three-digit numbers that do not have a digit $2$ or a digit $3$ is $\#(A \cup B) = 848$.

 There are $9 \cdot 10 \cdot 10 = 900$ ways to construct a three-digit number in total, so we have $900 - 848 = \boxed{\textbf{(A) }52}$ numbers that have at least one of digit $2$ and $3$.

 


	\item (\textit{AMC 12 12P 2002}) Let $AB$ be a segment of length $26$, and let points $C$ and $D$ be located on $AB$ such that $AC=1$ and $AD=8$. Let $E$ and $F$ be points on one of the semicircles with diameter $AB$ for which $EC$ and $FD$ are perpendicular to $AB$. Find $EF.$
$\text{(A) }5 \qquad \text{(B) }5 \sqrt{2}  \qquad \text{(C) }7 \qquad \text{(D) }7 \sqrt{2} \qquad \text{(E) }12$

 \textbf{Answer: }D



\textbf{Solution}

We can solve this with some simple coordinate geometry. Let $A$ be the origin and let $AB$ be located on the positive $x-$axis. The equation of semi-circle $AB$ is $(x-13)^2+y^2=13^2, y \geq 0.$ Since $E$ and $F$ are both perpendicular to $C$ and $D$ respectively, they must have the same $x -$ coordinate. Plugging in $1$ and $8$ into our semi-circle equation gives us $y=5$ and $y=12$ respectively. The distance formula on $(1, 5)$ and $(8, 12)$ gives us our answer of $\sqrt{(1-8)^2 + (5-12)^2}=\sqrt{2(7^2)}=\boxed{\textbf{(D) } 7\sqrt{2}}.$



	\item (\textit{AMC 12 12P 2002}) Two walls and the ceiling of a room meet at right angles at point $P.$  A fly is in the air one meter from one wall, eight meters from the other wall, and nine meters from point $P$. How many meters is the fly from the ceiling?


 $\text{(A) }\sqrt{13} \qquad \text{(B) }\sqrt{14} \qquad \text{(C) }\sqrt{15} \qquad \text{(D) }4 \qquad \text{(E) }\sqrt{17}$

 \textbf{Answer: }D



\textbf{Solution}

We can use the formula for the diagonal of a rectangular prism, or $d=\sqrt{a^2+b^2+c^2}$ The problem gives us $a=1, b=8,$ and $d=9.$ Solving gives us: $9=\sqrt{1^2 + 8^2 + c^2} \implies c^2=9^2-8^2-1^2 \implies c^2=16  \implies c=\boxed{\textbf{(D) } 4}.$



	\item (\textit{AMC 12 12P 2002}) Let $f_n (x) = \text{sin}^n x + \text{cos}^n x.$ For how many $x$ in $[0,\pi]$ is it true that


 
 \[6f_{4}(x)-4f_{6}(x)=2f_{2}(x)?\]
$\text{(A) }2 \qquad \text{(B) }4  \qquad \text{(C) }6 \qquad \text{(D) }8 \qquad \text{(E) more than }8$

 \textbf{Answer: }E



\textbf{Solution 1}

Divide by 2 on both sides to get 
 \[3f_{4}(x)-2f_{6}(x)=f_{2}(x)\]Substituting the definitions of $f_{2}(x)$, $f_{4}(x)$, and $f_{6}(x)$, we may rewrite the expression as
 \[3(\text{sin}^4{x} + \text{cos}^4{x}) - 2(\text{sin}^6{x} + \text{cos}^6{x}) = 1\]We now simplify each term separately using some algebraic manipulation and the Pythagorean identity.

 We can rewrite $3(\text{sin}^4 x + \text{cos}^4 x)$ as$3(\text{sin}^2 x + \text{cos}^2 x)^2 - 6\text{sin}^2 x \text{cos}^2 x$, which is equivalent to $3 - 6\text{sin}^2 x \text{cos}^2 x$.

 As for $2(\text{sin}^6 x + \text{cos}^6 x)$, we may factor it as $2(\text{sin}^2 x + \text{cos}^2 x)(\text{sin}^4 x + \text{cos}^4 x - \text{sin}^2 x \text{cos}^2 x)$ which can be rewritten as$2(\text{sin}^4 x + \text{cos}^4 x - \text{sin}^2 x \text{cos}^2 x)$, and then as$2(\text{sin}^2 x + \text{cos}^2 x)^2 - 6\text{sin}^2 x \text{cos}^2 x$, which is equivalent to $2 - 6\text{sin}^2 x \text{cos}^2 x$.

 Putting everything together, we have $(3 - 6\text{sin}^2 {x} \text{cos}^2 {x}) - (2 - 6\text{sin}^2 {x} \text{cos}^2 {x}) = 1$ or $1 = 1$. Therefore, the given equation $3f_{4}(x)-2f_{6}(x)=f_{2}(x)$ is true for all real $x$, meaning that there are certainly more than $8$ values of $x$ in the interval $[0,\pi]$ that satisfy the given equation, and so the answer is $\boxed{\text{(E) more than }8}$.

 


	\item (\textit{AMC 12 12P 2002}) Let $t_n = \frac{n(n+1)}{2}$ be the $n$th triangular number. Find


 
 \[\frac{1}{t_1} + \frac{1}{t_2} + \frac{1}{t_3} + ... + \frac{1}{t_{2002}}\]
$\text{(A) }\frac {4003}{2003} \qquad \text{(B) }\frac {2001}{1001} \qquad \text{(C) }\frac {4004}{2003} \qquad \text{(D) }\frac {4001}{2001} \qquad \text{(E) }2$

 \textbf{Answer: }C



\textbf{Solution 1}

We may write $\frac{1}{t_n}$ as $\frac{2}{n(n+1)}$ and do a partial fraction decomposition.Assume $\frac{2}{n(n+1)} = \frac{A_1}{n} + \frac{A_2}{n+1}$. 

 Multiplying both sides by $n(n+1)$ gives $2 = A_1(n+1) + A_2(n) = (A_1 + A_2)n + A_1$.

 Equating coefficients gives $A_1 = 2$ and $A_1 + A_2 = 0$, so $A_2 = -2$. Therefore, $\frac{2}{n(n+1)} = \frac{2}{n} - \frac{2}{n+1}$.

 Now $\frac{1}{t_1} + \frac{1}{t_2} + ... + \frac{1}{t_{2002}} = 2((\frac{1}{1} - \frac{1}{2}) + (\frac{1}{2} - \frac{1}{3}) + ... + (\frac{1}{2002} - \frac{1}{2003})) = 2(1 - \frac{1}{2003}) = \boxed{\textbf{(C) } \frac{4004}{2003}}.$

 Note: For the sake of completeness, I put the full derivation of the partial fraction decomposition of $\frac{2}{n(n+1)}$ here. However, on the contest, the decomposition step would be much faster since it is so well-known.

 



\textbf{Solution 2 (Cheese)}

As with all telescoping problems, there is a solution that involves induction. In competition, it is sufficient to conjecture the formula but not prove it. For sake of completeness and practice, we will prove the formula for $\sum_{i=1}^{n} \frac{1}{t_n}.$

 With some guess and check:

 $n=1$

 \begin{align*}\frac{1}{t_1}&=\frac{2}{1(2)} \\&=\frac{2}{2} \\\end{align*}


 $n=2$

 \begin{align*}\frac{1}{t_1}+\frac{1}{t_2}&=\frac{2}{2}+\frac{2}{2(3)} \\&=\frac{4}{3} \\\end{align*}


 $n=3$

 \begin{align*}\frac{1}{t_1}+\frac{1}{t_2}+\frac{1}{t_3}&=\frac{4}{3}+\frac{2}{3(4)} \\=\frac{6}{4} \\\end{align*}


 From $\frac{2}{2}, \frac{4}{3},$ and $\frac{6}{4},$ we can conjecture $\sum_{i=1}^{n} \frac{1}{t_n} = \frac{2n}{n+1}.$ A quick check shows that for $n=4$ gives $\frac{8}{5},$ which means our inductive hypothesis is most likely correct. Thus, our answer is $\sum_{i=1}^{2002} \frac{1}{t_n} = \frac{2(2002)}{2002+1}=\boxed{\textbf{(C) } \frac{4004}{2003}}.$ 

 We will prove this with induction for all $n \geq 1$ as promised.

 Base case: $n=1.$

 \begin{align*}\sum_{i=1}^{1} \frac{1}{t_n}&=\frac{1}{t_n} \\&=\frac{2}{1(2)} \\&=1 \\&=\frac{2(1)}{1+1} \\&=1 \\\end{align*}


 Since $1=1,$ the base case is proven.

 Induction step: $n=k+1.$

 Assume: $n=k$ $\implies$ $\sum_{i=1}^{k} \frac{1}{t_k} = \frac{2k}{k+1}.$

 We want to prove: $n=k+1$ $\implies$ $\sum_{i=1}^{k+1} \frac{1}{t_k} = \frac{2(k+1)}{k+2}.$

 We are given $\sum_{i=1}^{k} \frac{1}{t_k} = \frac{2k}{k+1}.$ Add $\frac{1}{t_{k+1}}$ on both sides and simplify.

 \begin{align*}\sum_{i=1}^{k} \frac{1}{t_k} + \frac{1}{t_{k+1}} &= \frac{2k}{k+1} + \frac{1}{t_{k+1}} \\\sum_{i=1}^{k+1} \frac{1}{t_k} &= \frac{2k}{k+1} + \frac{2}{(k+1)(k+2)} \\\sum_{i=1}^{k+1} \frac{1}{t_k} &= \frac{2k(k+2)}{(k+1)(k+2)} + \frac{2}{(k+1)(k+2)} \\\sum_{i=1}^{k+1} \frac{1}{t_k} &= \frac{2k(k+2)+2}{(k+1)(k+2)} \\\sum_{i=1}^{k+1} \frac{1}{t_k} &= \frac{2k^2+4k+2}{(k+1)(k+2)} \\\sum_{i=1}^{k+1} \frac{1}{t_k} &= \frac{2(k+1)^2}{(k+1)(k+2)} \\\sum_{i=1}^{k+1} \frac{1}{t_k} &= \frac{2(k+1)}{(k+2)} \\\end{align*}


 Given the inductive hypothesis, we have proven the induction step. Thus, we have completed our proof for induction.

 


	\item (\textit{AMC 12 12P 2002}) For how many positive integers $n$ is $n^3 - 8n^2 + 20n - 13$ a prime number?


 $\text{(A) one} \qquad \text{(B) two} \qquad \text{(C) three} \qquad \text{(D) four} \qquad \text{(E) more than four}$

 \textbf{Answer: }C



\textbf{Solution 1}

Since this is a number theory question, it is clear that the main challenge here is factoring the given cubic. In general, the rational root theorem will be very useful for these situations. 

 The rational root theorem states that all rational roots of $n^3 - 8n^2 + 20n - 13$ will be among $1, 13, -1$, and $-13$. Evaluating the cubic at these values will give $n = 1$ as a root. Doing some synthetic division gives 
 \[n^3 - 8n^2 + 20n - 13 = (n-1)(n^2 - 7n + 13)\]

 Since $n > 0$, $n-1$ must be positive. Since $(n-1)(n^2 - 7n + 13)$ evaluates to a prime, it is clear that exactly one of $n-1$ and $n^2 - 7n - 13$ is $1$. We proceed by splitting the problem into 2 cases.

 Case 1: $n-1 = 1$It is clear that $n = 2$.  We have $2^2 - 7(2) + 13 = 3$, so this case yields $n = 2$ as a solution.

 Case 2: $n^2 - 7n + 13 = 1$Solving for $n$ gives $n^2 - 7n + 12 = 0$ or $(n-3)(n-4) = 0$. Therefore, $n = 3$ or $n = 4$.Since both $3-1 = 2$ and $4-1 = 3$ are prime, both $n = 3$ and $n = 4$ work, yielding 2 solutions.

 Putting everything together, the answer is $1 + 2 = \boxed{\textbf{(C) }3}$.

 


	\item (\textit{AMC 12 12P 2002}) What is the maximum value of $n$ for which there is a set of distinct positive integers $k_1, k_2, ... k_n$ for which


 
 \[k^2_1 + k^2_2 + ... + k^2_n = 2002?\]
$\text{(A) }14 \qquad \text{(B) }15 \qquad \text{(C) }16 \qquad \text{(D) }17 \qquad \text{(E) }18$

 \textbf{Answer: }D



\textbf{Solution}

Note that $k^2_1 + k^2_2 + ... + k^2_n = 2002 \geq \frac{n(n+1)(2n+1)}{6}$

 When $n = 17$, $\frac{n(n+1)(2n+1)}{6} = \frac{(17)(18)(35)}{6} = 1785 < 2002$.

 When $n = 18$, $\frac{n(n+1)(2n+1)}{6} = 1785 + 18^2 = 2109 > 2002$.

 Therefore, we know $n \leq 17$.

 Now we must show that $n = 17$ works. We replace some integer $b$ within the set $\{1, 2, ... 17\}$ with an integer $a > 17$ to account for the amount under $2002$, which is $2002-1785 = 217$.

 Essentially, this boils down to writing $217$ as a difference of squares. Assume there exist positive integers $a$ and $b$ where $a > 17$ and $b \leq 17$ such that $a^2 - b^2 = 217$.

 We can rewrite this as $(a+b)(a-b) = 217$. Since $217 = 7 \cdot 31$, either $a+b = 217$



	\item (\textit{AMC 12 12P 2002}) Find $i + 2i^2 +3i^3 + . . . + 2002i^{2002}.$
$\text{(A) }-999 + 1002i \qquad \text{(B) }-1002 + 999i \qquad \text{(C) }-1001 + 1000i \qquad \text{(D) }-1002 + 1001i \qquad \text{(E) }i$

 \textbf{Answer: }D



\textbf{Solution 1}

Note that $i^4 = 1$, so $i^n = i^{4m+n}$ for all integers $m$ and $n$. In particular, $i = 1$, $i^2 = -1$, and $i^3 = -i$.We group the positive and negative real terms together and group the positive and negative imaginary parts together.

 The positive real terms have exponents on $i$ that are multiples of 4. Therefore, the positive real part evaluates to 
 \[4 + 8 + ... + 2000\]The negative real terms have exponents on $i$ that are of the form $4k + 2$ for integers $k$. Therefore, the negative real part evaluates to 
 \[-(2 + 6 + ... + 2002)\]The positive imaginary terms have exponents on $i$ that are of the form $4k + 1$ for integers $k$. Therefore, the negative real part evaluates to 
 \[(1 + 5 + ... + 2001)i\]The negative imaginary terms have exponents on $i$ that are of the form $4k + 3$ for integers $k$. Therefore, the negative real part evaluates to 
 \[-(3 + 7 + ... + 1999)i\]

 Putting everything together, we have 
 \[i + 2i^2 + ... + 2002i^{2002} = (-2 + 4 - 6 + 8 - ... + 2000 - 2002) + (1 - 3 + 5 - 7 + ... - 1999 + 2001)i\]

 Group every 2 consecutive terms as shown below 
 \[((-2 + 4)+(-6 + 8) + ... + (-1998 + 2000)-2002) + ((1 - 3)+(5 - 7) + ... + (1997 - 1999) + 2001)i\]

 Now we evaluate each small bracket with 2 terms. We get $500(2) = 1000$ in the real part and $500(-2) = -1000$ in the imaginary part. Therefore, the sum becomes $(1000 - 2002) + (-1000 + 2001)i = \boxed {\text{(D) }-1002 + 1001i}$.

 Note: This problem is similar to \href{https://artofproblemsolving.com/wiki/index.php?title=2009\_AMC\_12A\_Problems/Problem\_15}{2009 AMC 12A Problem 15}

 


	\item (\textit{AMC 12 12P 2002}) There are $1001$ red marbles and $1001$ black marbles in a box. Let $P_s$ be the probability that two marbles drawn at random from the box are the same color, and let $P_d$ be the probability that they are different colors. Find $|P_s-P_d|.$
$\text{(A) }0 \qquad \text{(B) }\frac{1}{2002} \qquad \text{(C) }\frac{1}{2001} \qquad \text{(D) }\frac {2}{2001} \qquad \text{(E) }\frac{1}{1000}$

 \textbf{Answer: }C



\textbf{Solution}

First we find the value of $P_s$. Note that whatever color we choose on our first marble, there are exactly $1000$ of $2001$ marbles remaining that match that color. Therefore, $P_s = \frac {1000}{2001}$.

 Now we find the value of $P_d$. Again, the actual color of the first marble does not matter, since there are always exactly $1001$ of $2001$ marbles remaining that don't match that color. Therefore, $P_d = \frac{1001}{2001}$.

 The value of $|P_s - P_d|$ is therefore $\frac {|1000-1001|}{2001} = \boxed {\text{(C) }\frac{1}{2001}}$.

 


	\item (\textit{AMC 12 12P 2002}) The altitudes of a triangle are $12, 15,$ and $20.$ The largest angle in this triangle is


 $\text{(A) }72^\circ \qquad \text{(B) }75^\circ \qquad \text{(C) }90^\circ \qquad \text{(D) }108^\circ \qquad \text{(E) }120^\circ$

 \textbf{Answer: }C



\textbf{Solution}

Let $a, b,$ and $c$ denote the bases of altitudes $12, 15,$ and $20,$ respectively. Since they are all altitudes and bases of the same triangle, they have the same area, so $\frac{12a}{2}=\frac{15b}{2}=\frac{20c}{2}.$ Multiplying by $2$, we get $12a=15b=20c.$ Notice that a simple solution to the equation is if all of them equal $12 \cdot 15 \cdot 20.$ That means $a=15 \cdot 20, b=12 \cdot 20,$ and $c=12 \cdot 15.$ Simplifying our solution to check for Pythagorean triples we see that this is just a Pythagorean triple, namely a $3-4-5$ triangle. Since the other two angles of a right triangle must be acute, the right angle must be the greatest angle. Therefore, our answer is $\boxed{\textbf{(C) }90^\circ}.$



	\item (\textit{AMC 12 12P 2002}) Let $f(x) = \sqrt{\sin^4{x} + 4 \cos^2{x}} - \sqrt{\cos^4{x} + 4 \sin^2{x}}.$ An equivalent form of $f(x)$ is


 $\text{(A) }1-\sqrt{2}\sin{x} \qquad \text{(B) }-1+\sqrt{2}\cos{x} \qquad \text{(C) }\cos{\frac{x}{2}} - \sin{\frac{x}{2}} \qquad \text{(D) }\cos{x} - \sin{x} \qquad \text{(E) }\cos{2x}$

 \textbf{Answer: }E



\textbf{Solution 1}

By the Pythagorean identity we can rewrite the given expression as follows.
 \[\sqrt{\sin^4{x} + 4 \cos^2{x}} - \sqrt{\cos^4{x} + 4 \sin^2{x}} = \sqrt{\sin^4{x} + 4(1 - \sin^2{x})} - \sqrt{\cos^4{x} + 4(1 - \cos^2{x})}\]

 Expanding each bracket gives 
 \[\sqrt{\sin^4{x} - 4\sin^2{x} + 4} - \sqrt{\cos^4{x} - 4\cos^2{x} + 4}\]

 The expressions under the square roots can be factored to get 
 \[\sqrt{(\sin^2{x} - 2)^2} - \sqrt{(\cos^2{x} - 2)^2}\]

 Since $\sin^2{x} < 2$ and $\cos^2{x} < 2$ for all real $x$, the expression must evaluate to $(2 - \sin^2{x}) - (2 - \cos^2{x})$, which simplifies to $\cos^2{x} - \sin^2{x} = \boxed {\text{(E) }\cos{2x}}$.

 



\textbf{Solution 2 (Cheese)}

We don't actually have to solve the question. Just let $x$ equal some easy value to calculate $\cos {x}, \cos {2x}, \sin {x}, \sin {\frac{x}{2}},$ and $\cos {\frac{x}{2}}.$ For this solution, let $x=60^\circ.$ This means that the expression in the problem will give $\sqrt{\sin^4{60^\circ} + 4 \cos^2{60^\circ}} - \sqrt{\cos^4{60^\circ} + 4 \sin^2{60^\circ}}=\sqrt{(\frac{\sqrt{3}}{2})^4 + 4(\frac{1}{2})^2} - \sqrt{(\frac{1}{2})^4 + 4(\frac{\sqrt{3}}{2})^2}=\sqrt{\frac{9}{16} +1} - \sqrt{\frac{1}{16} + 3} = \frac{-1}{2}.$ Plugging in $x=60^\circ$ for the rest of the expressions, we get

 $\text{(A) }1-\sqrt{2}\sin{60^\circ}=1-\frac{\sqrt{6}}{2} \neq \frac{-1}{2}.$
$\text{(B) }1+\sqrt{2}\cos{60^\circ}=1+\frac{\sqrt{2}}{2} \neq \frac{-1}{2}.$
$\text{(C) }\cos{\frac{60^\circ}{2}} - \sin{\frac{60^\circ}{2}}=\frac{\sqrt{3}}{2}-1 \neq \frac{-1}{2}.$
$\text{(D) }\cos{60^\circ} - \sin{60^\circ}=\frac{1-\sqrt{3}}{2} \neq \frac{-1}{2}.$
$\text{(E) }\cos {120^\circ}=\frac{-1}{2}.$

 Therefore, our answer is $\boxed{\textbf{(E) } \cos {2x}}$.

 Comment: If you decide to cheese the problem, be very careful not to choose any $x$ where $x$ is a multiple of $\frac{\pi}{4}$. It turns out that all of them except for $\frac{7\pi}{4} + 2\pi k$ result in equality between the correct answer and one or more wrong answers, which one could quickly verify by setting two choices equal and solving for $x$.

 


	\item (\textit{AMC 12 12P 2002}) If $a,b,c$ are real numbers such that $a^2 + 2b =7$, $b^2 + 4c= -7,$ and $c^2 + 6a= -14$, find $a^2 + b^2 + c^2.$
$\text{(A) }14 \qquad \text{(B) }21 \qquad \text{(C) }28 \qquad \text{(D) }35 \qquad \text{(E) }49$

 \textbf{Answer: }A



\textbf{Solution 1}

Adding all of the equations gives $a^2 + b^2 +c^2 + 6a + 2b + 4c=-14.$ Adding 14 on both sides gives $a^2 + b^2 +c^2 + 6a + 2b + 4c+14=0.$ Notice that 14 can split into $9, 1,$ and $4,$ which coincidentally makes $a^2 +6a, b^2+2b,$ and $c^2+4c$ into perfect squares. Therefore, $(a+3)^2 + (b+1)^2 + (c+2) ^2 =0.$ An easy solution to this equation is $a=-3, b=-1,$ and $c=-2.$ Plugging in that solution, we get $a^2+b^2+c^2=(-3)^2+(-1)^2+(-2)^2=\boxed{\textbf{(A) } 14}.$



	\item (\textit{AMC 12 12P 2002}) In quadrilateral $ABCD$, $m\angle B = m \angle C = 120^{\circ}, AB=3, BC=4,$ and $CD=5.$ Find the area of $ABCD.$
$\text{(A) }15 \qquad \text{(B) }9 \sqrt{3} \qquad \text{(C) }\frac{45 \sqrt{3}}{4} \qquad \text{(D) }\frac{47 \sqrt{3}}{4} \qquad \text{(E) }15 \sqrt{3}$

 \textbf{Answer: }D



\textbf{Solution 1}

Draw $AE$ parallel to $BC$ and draw $BF$ and $CG \perp AE$, where $F$ and $G$ are on $AE$.

 It is clear that triangles $AFB$ and $EGC$ are congruent $30-60-90$ triangles. Therefore, $AF = EG = \frac{3}{2}$ and $BF = CG = \frac{3\sqrt{3}}{2}$.

 Therefore, $AE = AF + FG + GC = 4 + (2)(\frac{3}{2}) = 7$ and the area of trapezoid $ABCE$ is $(\frac{1}{2})(4+7)(\frac{3\sqrt{3}}{2}) = \frac{33\sqrt{3}}{4}$.

 It remains to find the area of triangle $AED$, which is $(\frac{1}{2})(AE)(ED)(\sin 120^{\circ}) = (\frac{1}{2})(7)(2)(\frac{\sqrt{3}}{2}) = \frac{7\sqrt{3}}{2}$.

 Therefore, the total area of quadrilateral $ABCD$ is $\frac{33\sqrt{3}}{4} + \frac{7\sqrt{3}}{2} = \boxed{\textbf{(D) }\frac{47\sqrt{3}}{4}}$.

 



\textbf{Solution 2}

Extend $AB$ and $CD$ to meet at $E$. Then triangle $BEC$ is equilateral, as $\angle EBC = \angle ECB = 60^{\circ}$. Thus, $BC = CE = EB = 4$. 

 Note $AE = AB + BE = 7$ and $DE = DC + CE = 9$ Then, the area of triangle $ADE$ is $\frac {1} {2} \cdot AE \cdot DE \cdot \sin{60^{\circ}} = \frac {63\sqrt{3}} {4}$. As triangle $BEC$ is equilateral, its area is $\frac {\sqrt{3}} {4} \cdot 4^2 = \frac {16\sqrt{3}} {4}$. 

 The area of $ABCD$ is the area of triangle $ADE$ minus the area of triangle $BEC$, or $\boxed{\textbf{(D) }\frac{47\sqrt{3}}{4}}$.

 


	\item (\textit{AMC 12 12P 2002}) Let $f$ be a real-valued function such that


 
 \[f(x) + 2f(\frac{2002}{x}) = 3x\]


 for all $x>0.$ Find $f(2).$
$\text{(A) }1000 \qquad \text{(B) }2000 \qquad \text{(C) }3000 \qquad \text{(D) }4000 \qquad \text{(E) }6000$

 \textbf{Answer: }B



\textbf{Solution}

Setting $x = 2$ gives $f(2) + 2f(1001) = 6$.Setting $x = 1001$ gives $2f(2) + f(1001) = 3003$.

 Adding these 2 equations and dividing by 3 gives$f(2) + f(1001) = \frac{6+3003}{3} = 1003$.

 Subtracting these 2 equations gives$f(2) - f(1001) = 2997$.

 Therefore, $f(2) = \frac{1003+2997}{2} = \boxed {\textbf{(B) }2000}$.

 \textbf{Note} We can find a closed form expression for $f$ using this method. We have that 
 \[f(x) + 2f(\frac{2002}{x}) = 3x\] and plugging in $\frac{2002}{x}$ gives 
 \[f(\frac{2002}{x}) + 2f(x) = 3(\frac{2002}{x}).\]

 Let $a = f(x)$ and $b = f(\frac{2002}{x})$. Then, we have 
 \[a + 2b = 3x \tag{1}\] and 
 \[2a + b = \frac{6006}{x} \tag{2}\]

 We add equations $(1)$ and $(2)$ and divide by $3$ to get $a+b = x + \frac{2002}{x}$. Subtracting this from equation $(2)$ gives us the closed form 
 \[f(x) = \frac{4004}{x} - x\]

 Indeed, plugging in $x=2$ gives us $f(2) = \frac{4004}{2} - 2 = \boxed {\textbf{(B) }2000}$.

 



\textbf{Solution 2(Similar to Sol 1)}

Setting $x = 2$ in the given equation yields
 \[f(2) + 2f(1001) = 6.\]Setting $x = 1001$ gives
 \[2f(2) + f(1001) = 3003.\]

 Adding these two equations and dividing by $3$ results in
 \[f(2) + f(1001) = \frac{6 + 3003}{3} = 1003.\]

 We now have the system:\begin{align*}f(2) + f(1001) &= 1003 \\2f(2) + f(1001) &= 3003\end{align*}


 Subtracting the first equation from the second gives
 \[f(2) = 3003 - 1003 = 2000.\]

 Thus, $f(2) = \boxed{\textbf{(B) }2000}$.

 


	\item (\textit{AMC 12 12P 2002}) Let $a$ and $b$ be real numbers greater than $1$ for which there exists a positive real number $c,$ different from $1$, such that


 
 \[2(\log_a{c} + \log_b{c}) = 9\log_{ab}{c}.\]


 Find the largest possible value of $\log_a b.$
$\text{(A) }\sqrt{2} \qquad \text{(B) }\sqrt{3} \qquad \text{(C) }2 \qquad \text{(D) }\sqrt{6} \qquad \text{(E) }3$

 \textbf{Answer: }C



\textbf{Solution}

We may rewrite the given equation as 
 \[2 \left(\frac {\log c}{\log a} + \frac {\log c}{\log b} \right) = \frac {9\log c}{\log a + \log b}\]Since $c \neq 1$, we have $\log c \neq 0$, so we may divide by $\log c$ on both sides. After making the substitutions $x = \log a$ and $y = \log b$, our equation becomes 
 \[\frac {2}{x} + \frac {2}{y} = \frac {9}{x+y}\]

 Rewriting the left-hand side gives 
 \[\frac {2(x+y)}{xy} = \frac {9}{x+y}\]

 Cross-multiplying gives $2(x+y)^2 = 9xy$ or 
 \[2x^2 - 5xy + 2y^2 = 0\]

 Factoring gives $(2x-y)(x-2y) = 0$ or $\frac {x}{y} = 2, \frac {1}{2}$.

 Recall that $\frac {x}{y} = \frac {\log a}{\log b} = \log_{a} b$. Therefore, the maximum value of $\log_{a} b$ is $\boxed {\text{(C) }2}$.

 


	\item (\textit{AMC 12 12P 2002}) Under the new AMC $10, 12$ scoring method, $6$ points are given for each correct answer, $2.5$ points are given for each unanswered question, and no points are given for an incorrect answer. Some of the possible scores between $0$ and $150$ can be obtained in only one way, for example, the only way to obtain a score of $146.5$ is to have $24$ correct answers and one unanswered question. Some scores can be obtained in exactly two ways, for example, a score of $104.5$ can be obtained with $17$ correct answers, $1$ unanswered question, and $7$ incorrect, and also with $12$ correct answers and $13$ unanswered questions. There are (three) scores that can be obtained in exactly three ways. What is their sum?


 $\text{(A) }175 \qquad \text{(B) }179.5 \qquad \text{(C) }182 \qquad \text{(D) }188.5 \qquad \text{(E) }201$

 \textbf{Answer: }D



\textbf{Solution 1}

Let $c$ denote the amount of correct answers, $u$ denote the amount of unanswered questions, and $w$ denote the amount of wrong answers. The conversion rate is what differentiates between two ways to get the same score.

 There are two ways to realize the conversion rate of wrong answers to unanswered questions and correct answers to unanswered questions.

 \textbf{Sub Solution 1.1} Notice how the way to get $104.5$ points directly tells us the conversion rate. $17$ correct questions, $1$ unanswered question, and $7$ wrong answers is equivalent to $12$ correct answers, $13$ unanswered questions, and $0$ wrong answers. Since $17-12=5, 1-13=-12,$ and $7-0=7$, the conversion rate is $c, u, w \implies c-5, u+12, w-7.$

 \textbf{Sub Solution 1.2} The second way to find out the conversion rate is to actually do the math. If a wrong answer is converted to an unanswered question, the total value will be $2.5$ more than the total value. If a correct answer is converted to an unanswered question, the total value will be $6-2.5=3.5$ less than the total value. Thus, we want the "more than" to be equal to the "less than". 

 Let $x$ denote the amount of correct answers converted to unanswered questions and $y$ denote the amount of wrong answers converted to unanswered questions. Since we want "more than" to be equal to the "less than", $3.5x=2.5y$, where $x$ and $y$ are integers. Multiply by $2$ on both sides, $7x=5y$ implies $x=5$ and $y=7.$ Again, the conversion rate is $c, u, w \implies c-5, u+12, w-7.$ 

 \textbf{Sub Solutions Combine} Either way, we find $c, u, w \implies c-5, u+12, w-7.$ If there are three ways, this implies the first way $c, u, w$ the second way is $c, u, w \implies c-5, u+12, w-7,$ and the third way is $c, u, w \implies c-10, u+24, w-14.$ Since $c, u, w \geq0,$ it follows that our three cases' first ways are as follows:

 
$11$ correct answers, $0$ unanswered questions, $14$ wrong answers

 $10$ correct answers, $0$ unanswered questions, $15$ wrong answers

 $10$ correct answers, $1$ unanswered questions, $14$ wrong answers.

 
Thus, our answer is $6(11+10+10)+2.5(0+0+1)+0(14+15+14)=\boxed {\text{(D) }188.5}.$




\textbf{Solution 2}

This solution is formulated in formal language rather than cheesy methods, though the idea is the same. Additionally, we do not need to take into account the amount of wrong answers, which is easier to deal with.

 Let $c$ denote the amount of correct answers, $u$ denote the amount of unanswered questions. Then $c+u \leq 25.$ We want to find the LCM of $2.5$ and $6$ so we can convert between different ways, similar to solution $1.$ The LCM is equal to $30$, since $6 \cdot \textbf{5}=30$ and $2.5 \cdot \textbf{12}=30.$ 

 The first way will take the form $(c,u).$ The second way will take the form $(c+5, u-12).$ The third way will take the form $(c+10, u-24).$ In particular, since we cannot have a negative amount of unanswered questions, $u \geq 24.$ Combined with $c+u \leq 25$, our three cases are $(0,24), (1,24),$ and $(0,25).$

 Thus, our answer is $6(0+1+0) + 2.5(24+24+25)=\boxed {\text{(D) }188.5}.$

 \textbf{Note} The conversion rate can either be $(c+5, u-12).$ or $(c-5, u+12).$ Both solutions use a different conversion rate to illustrate this point. The first solution defines the "first way" as the one with the most correct, while the second solution defines the "first way" as the one with the most unanswered.

 


	\item (\textit{AMC 12 12P 2002}) The equation $z(z+i)(z+3i)=2002i$ has a zero of the form $a+bi$, where $a$ and $b$ are positive real numbers. Find $a.$
$\text{(A) }\sqrt{118} \qquad \text{(B) }\sqrt{210} \qquad \text{(C) }2 \sqrt{210} \qquad \text{(D) }\sqrt{2002} \qquad \text{(E) }100 \sqrt{2}$

 \textbf{Answer: }A



\textbf{Solution}

Note that $2002 = 11 \cdot 13 \cdot 14$. With this observation, it becomes easy to note that $z = -14i$ is a root of the given equation. However, it is not of the desired form in the problem, so we must factor the given expression to obtain the other 2 roots. From this point onwards, we assume that $z \neq -14i$.

 Expanding $z(z+i)(z+3i)=2002i$, we have $z^3 + 4iz^2 - 3z - 2002i = 0$. We may factor it as $(z + 14i)(z^2 - 10iz - 143) = 0$. Since $z + 14i \neq 0$, we must have $z^2 - 10iz - 143 = 0$. Therefore, $z = \frac {10i \pm \sqrt{-100 + 4(143)}}{2} = 5i \pm \sqrt {118}$.

 Since $a > 0$, we ignore the negative root. Therefore, $a = \boxed {\text{(A) } \sqrt {118}}$.

 


	\item (\textit{AMC 12 12P 2002}) Let $ABCD$ be a regular tetrahedron and Let $E$ be a point inside the face $ABC.$ Denote by $s$ the sum of the distances from $E$ to the faces $DAB, DBC, DCA,$ and by $S$ the sum of the distances from $E$ to the edges $AB, BC, CA.$ Then $\frac{s}{S}$ equals


 $\text{(A) }\sqrt{2} \qquad \text{(B) }\frac{2 \sqrt{2}}{3} \qquad \text{(C) }\frac{\sqrt{6}}{2} \qquad \text{(D) }2 \qquad \text{(E) }3$

 \textbf{Answer: }B



\textbf{Solution}

Assume points $P$, $Q$, and $R$ are on faces $ABD$, $ACD$, and $BCD$ respectively such that $EP \perp ABD$, $EQ \perp ACD$, and $ER \perp BCD$.

 Assume points $S$, $T$, and $U$ are on edges $AB$, $AC$, and $BC$ respectively such that $ES \perp AB$, $ET \perp AC$, and $EU \perp BC$.

 Consider triangles $EPS$, $EQT$, and $ERU$. Each of these triangles have a right angle and an angle equal to the dihedral angle of the tetrahedron, so they are all similar by AA similarity. In particular, we know that $\frac{EP}{ES} = \frac{EQ}{ET} = \frac{ER}{EU} = \frac{EP+EQ+ER}{ES+ET+EU} = \frac{s}{S}$.

 It remains to find $\frac{EP}{ES}$, or equivalently, $\sin(\angle DSE)$.

 We know $SE = \frac{1}{3}DS$ by the centroid property. Therefore, $\cos(\angle DSE) = \frac{1}{3}$, so $\sin(\angle DSE) = \sqrt{1-\left(\frac{1}{3}\right)^2} = \boxed {\text{(B) }\frac{2 \sqrt{2}}{3}}$.

 



\textbf{Solution 2(Cheesy)}

Continue to assume points $P$, $Q$, and $R$ are on faces $ABD$, $ACD$, and $BCD$ respectively such that $EP \perp ABD$, $EQ \perp ACD$, and $ER \perp BCD$ and assume points $S$, $T$, and $U$ are on edges $AB$, $AC$, and $BC$ respectively such that $ES \perp AB$, $ET \perp AC$, and $EU \perp BC$. 

 Now, because they never specify where E has to be or how long the tetrahedron's sides must be, WLOG, assume E is the centroid/incenter/circumcenter/orthocenter of triangle ABC and the side length of the tetrahedron is $2$ Note that the inradius of ABC is the same as $ES=ET=EU.$ Then the inradius of $\triangle\text{ABC}$ are $\frac{\sqrt3}{6}\cdot 2=\frac{\sqrt3}{3}.$ 

 Next, take cross section $\triangle\text{SDC}.$ Since $ES$ is an inradius, $ES=\frac{\sqrt3}{3}.$ Since $SD$ is an altitude of $\triangle\text{DAB}$ and $SC$ is an altitude of $\triangle\text{ABC}$, $SC=SD=\frac{\sqrt3}{2}\cdot 2=\sqrt3.$ Thus, by the Pythagorean theorem, $ED=\sqrt{SD^2-ES^2}=\sqrt{\sqrt{3}^2-\frac{\sqrt{3}^2}{3^2}}=\sqrt{3-\frac13}=\sqrt{\frac83}=\frac{2\sqrt6}{3}.$ 

 Then, because $\angle{SED}=\angle{SQE}=90^\circ$ and $\angle{ESD}=\angle{QSE},$ $\triangle{SED}~\triangle{SQE}$ by AA Similarity. Then, since $ES=\frac{SD}3=\frac{\sqrt3}3,$ $EQ=\frac{ED}3=\frac{\frac{2\sqrt6}3}{3}=\frac{2\sqrt6}9.$ 

 Then, $s=3\cdot EQ=3\cdot\frac{2\sqrt6}9=\frac{2\sqrt6}3$ and $S=3\cdot\text{inradius}=3\cdot\frac{\sqrt3}3=\sqrt3.$ Thus, $\frac{s}{S}=\frac{\frac{2\sqrt6}3}{\sqrt3}=\frac{2\sqrt2}3=\boxed{\text{(B)}}.$



	\item (\textit{AMC 12 12P 2002}) Let $a$ and $b$ be real numbers such that $\sin{a} + \sin{b} = \frac{\sqrt{2}}{2}$ and $\cos {a} + \cos {b} = \frac{\sqrt{6}}{2}.$ Find $\sin{(a+b)}.$
$\text{(A) }\frac{1}{2} \qquad \text{(B) }\frac{\sqrt{2}}{2} \qquad \text{(C) }\frac{\sqrt{3}}{2} \qquad \text{(D) }\frac{\sqrt{6}}{2} \qquad \text{(E) }1$

 \textbf{Answer: }C



\textbf{Solution}

Sum to product gives us

 
 \[2\sin(\frac{a+b}{2})\cos(\frac{a-b}{2}) = \frac{\sqrt{2}}{2}\]

 \[2\cos(\frac{a+b}{2})\cos(\frac{a-b}{2}) = \frac{\sqrt{6}}{2}\]

 Dividing these equations tells us that $\tan(\frac{a+b}{2}) = \frac{1}{\sqrt{3}}$, so $\frac{a+b}{2} = \frac{\pi}{6} + \pi n$ for an integer $n$, so $\sin(a+b) = \sin (\frac{\pi}{3} + 2\pi n) = \frac{\sqrt{3}}{2}$. The answer is $\boxed{\text{(C)}}$.

 


\end{enumerate}